\subsection{Complex Numbers} 

%\DTMnewdatestyle{dateupdated}{〈definition 〉

Answers to \thesubsection ~problems are on page \pageref{complex_numbers_answers}.  \textbf{\textcolor{red}{Updated \DTMnow.}}

%--------------------------------------------------------------------------------------------------------------
\begin{Exercise}
Find the real and imaginary parts of $4i - 3$.
\end{Exercise}
\begin{Answer}
4 and $-3$, respectively
\end{Answer}

%--------------------------------------------------------------------------------------------------------------
\begin{Exercise}
Find the magnitude (aka ``modulus'' or ``absolute value'') of $4-3i$.
\end{Exercise}
\begin{Answer}
5
\end{Answer}

%--------------------------------------------------------------------------------------------------------------
\begin{Exercise}
Find the phase (aka ``argument'') of $4+3i$.
\end{Exercise}
\begin{Answer}
0.643
\end{Answer}

%--------------------------------------------------------------------------------------------------------------
\begin{Exercise}
Find the phase of (a) $-4+3i$ and (b) $-4-3i$.  (Hint: your calculator lies.  Graph it.)
\end{Exercise}
\begin{Answer}
(a) 2.50 (b) 3.78 or $-2.50$
\end{Answer}

%--------------------------------------------------------------------------------------------------------------
\begin{Exercise}
Find the real and imaginary parts of $6e^{i \pi/6}$.
\end{Exercise}
\begin{Answer}
5.20 and 3, respectively
\end{Answer}

%--------------------------------------------------------------------------------------------------------------
\begin{Exercise}
Find the magnitude and phase of $6e^{i(4\pi /7)}$.
\end{Exercise}
\begin{Answer}
6 and $4\pi / 7$
\end{Answer}

%--------------------------------------------------------------------------------------------------------------
\begin{Exercise}
Find the real and imaginary parts of $(3+2i)(-4-6i)$.
\end{Exercise}
\begin{Answer}
0 and $-26$
\end{Answer}

%--------------------------------------------------------------------------------------------------------------
\begin{Exercise}
Find the real and imaginary parts of $(2+3i)/(4-5i)$.  (Hint: multiply both the numerator and denominator by the denominator's complex conjugate.)
\end{Exercise}
\begin{Answer}
$-0.1707$ and 0.537
\end{Answer}

%--------------------------------------------------------------------------------------------------------------
\begin{Exercise}
Plot each of the following on the complex plane: $e^{i\cdot 0}$, $e^{i\pi /2}$, $e^{i\pi}$, and $e^{i ∙3\pi / 2}$.
\end{Exercise}
\begin{Answer}
Your points are the same as $1, i, -1$, and $-i$.
\end{Answer}

%--------------------------------------------------------------------------------------------------------------
\begin{Exercise}
Write $4+3i$ in the form $re^{i\theta}$.
\end{Exercise}
\begin{Answer}
$5e^{0.643i}$
\end{Answer}

%--------------------------------------------------------------------------------------------------------------
\begin{Exercise}
What are the magnitude and phase of $e^{i\pi/6} e^{i\pi/4}$?
\end{Exercise}
\begin{Answer}
1 and $5\pi/12$
\end{Answer}

%--------------------------------------------------------------------------------------------------------------
\begin{Exercise}
What are the magnitude and phase of $(4+3i) e^{0.2i}$?
\end{Exercise}
\begin{Answer}
5 and 0.843
\end{Answer}

%--------------------------------------------------------------------------------------------------------------
\begin{Exercise}
Suppose $f(t)= 20e^{i\omega t}$, where $\omega=5$~sec$^{-1}$.  Find the magnitude and phase of $f(t)$ at (a)~$t=0$~seconds and (b)~at $t=0.27$~seconds.
\end{Exercise}
\begin{Answer}
(a) 20 and 0 (b) 20 and 1.35
\end{Answer}

%--------------------------------------------------------------------------------------------------------------
\begin{Exercise}
Suppose $f(t)= (6+8i) e^{i\omega t}$, where $\omega=5$~sec$^{-1}$.  Find the magnitude and phase of $f(t)$ at (a)~$t=0$~seconds and (b)~at $t=0.27$~seconds.
\end{Exercise}
\begin{Answer}
(a) 10 and 0.927 (b) 10 and 2.277 
\end{Answer}

%--------------------------------------------------------------------------------------------------------------
\begin{Exercise}
The ``exponential rule'' for derivatives of exponential functions is that $\displaystyle \frac{d}{dx} e^u=e^u  \frac{du}{dx}$.  
Show by direct substitution of the Euler formula, $e^{i \theta}=\cos \theta + i \sin \theta$, that the exponential rule yields the correct results for  $\displaystyle \frac{d}{dt} (Ae^{i\omega t}) \text{~and~}  \frac{d}{dt}(Ae^{-i\omega t})$.
\end{Exercise}




%\newcommand{\figuretoright}{3} {%
	%#1 is the picture width
	%#2 is the text to the right



%Ideas for additional problems:

%\vfill

\vspace{2in}

\pagebreak[3]

\textbf{Answers to Selected {\thesubsection} Problems:}
\label{complex_numbers_answers}
%\shipoutExercise
\shipoutAnswer

\cleardoublepage
